\documentclass[11pt,a4paper]{article}
\usepackage[T1]{fontenc}
\usepackage{isabelle,isabellesym}
\usepackage{pdfsetup}
\hypersetup{hidelinks}

\urlstyle{rm}
\isabellestyle{it}
\newcommand{\theoremname}[1]{{\urlstyle{tt}\nolinkurl{#1}}}
\AtBeginDocument{%
  \typeout{AFP-QA-TEXT-LEFT-SP=\number\dimexpr1in+\hoffset+\oddsidemargin\relax}%
  \typeout{AFP-QA-TEXT-RIGHT-SP=\number\dimexpr1in+\hoffset+\oddsidemargin+\textwidth\relax}%
}

\begin{document}
\pagenumbering{roman}

\title{Halting Absoluteness and Effective Sentence Invariance in Isabelle/ZF}
\author{Tang Ziyi}
\maketitle

\begin{abstract}
We construct a first-order set-theoretic formula whose satisfaction in every
transitive ZFC set model agrees with finite execution of a given set-coded
Turing machine. We then close the formula by naming its machine and input.
Let $I$ be the natural-number codes of closed formulas whose truth value
agrees across all such models. Assuming a countable transitive model of
ZFC exists, we construct primitive-recursive many-one reductions from both
self-input halting and its complement to $I$. The reductions use an
injective numerical code for formulas and the AFP formalisation of the
independence of the continuum hypothesis.
\end{abstract}

\tableofcontents
\clearpage
\pagenumbering{arabic}

\section{Halting formula}
For every machine $P$, input $x$, and transitive set model $M$ of ZFC, the
explicit formula $\mathit{halt\_fm}$ satisfies
\[
  (M,[P,x])\models\mathit{halt\_fm}
  \quad\Longleftrightarrow\quad
  \mathrm{halts}(P,x).
\]
Theorem \theoremname{transitive_model_halt_fm_iff} proves this statement.
The general method of absoluteness for transitive models is classical
\cite[Chapter~12]{jech}. Here the satisfaction theorem concerns the explicit
formula for this machine semantics. Finite run witnesses and their closure
in transitive models are established before the formula is introduced. The
closed sentence $\mathit{halt\_sentence}(P,x)$ names both parameters
internally. Its satisfaction equation is established by
\begin{quote}
\small\theoremname{sats_halt_sentence_iff}
\end{quote}
These results do not use CH or an assumption that a countable transitive
model exists.

\section{Effective reduction}
Let $K$ be the self-input halting set of the set-coded machines defined here,
and let $I$ contain exactly the numerical codes of closed formulas whose
truth value agrees across all transitive set models of ZFC. Given a
countable transitive set model of ZFC, this entry constructs
primitive-recursive $f_+$ and $f_-$ such that, for every natural number $e$,
\[
  e\in K\ \Longleftrightarrow\ f_+(e)\in I,
  \qquad
  e\notin K\ \Longleftrightarrow\ f_-(e)\in I.
\]
The formal reduction statements are
\begin{quote}
\small\ttfamily
self\_halting\_pr\_many\_one\_invariance\\
nonhalting\_pr\_many\_one\_invariance
\end{quote}
The unary notation means
$f_+(e)=\texttt{pr\_invariance\_positive}\mathbin{\texttt{`}}[e]$ and
$f_-(e)=\texttt{pr\_invariance\_negative}\mathbin{\texttt{`}}[e]$.
The formal primitive-recursive functions act on lists of natural numbers.
CH supplies a fixed sentence on which two transitive models
disagree, using the AFP formalisation of its independence~\cite{afp-ch}.
The reductions also hold under the direct premise that CH is not invariant.
The countable-model assumption supplies that premise through the
disagreement witnesses.

The target set contains only codes of well-formed closed formulas. A
constructor-tagged pairing code is proved injective on formulas. The two
reducing functions are proved members of Isabelle/ZF's class
\texttt{prim\_rec}, with application equations identifying their values
with the codes of the intended sentences. Their finite-list component uses
a bounded reverse traversal of the machine code.

The machine semantics follow the five-action design of the Isabelle/HOL
AFP entry \emph{Universal Turing Machine}~\cite{afp-utm}. Here the objects
and proofs are developed in Isabelle/ZF. Model closure, relativization, and
formula satisfaction use the AFP entry \emph{Transitive Models of Fragments
of ZFC}~\cite{afp-models}. A concrete rejection transformation
proves that $K$ is not decidable by a machine of this model. The two reductions
are certified in Isabelle/ZF's class \texttt{prim\_rec}. No theorem identifying
every member of that class with a machine of this model is assumed or proved.
Consequently, this entry does not claim a machine-decidability theorem for
the full target set $I$.

The set-theoretic formula datatype and its satisfaction relation are part
of Isabelle/ZF's constructibility development~\cite{paulson-zf}. The
primitive-recursive function class is supplied by the bundled
\texttt{ZF-Induct.Primrec} theory~\cite{szasz}.

\input{session}

\clearpage
\section*{Development provenance}
Tang Ziyi formulated the research question and wrote the initial
formalisation. Subsequent definitions, proofs, and repository engineering
include substantial assistance from OpenAI Codex. Aristotle-generated
draft proofs and review reports were also supplied during development.

\begin{thebibliography}{9}

\bibitem{afp-ch}
Emmanuel Gunther, Miguel Pagano, Pedro S\'anchez Terraf, and Mat\'ias Steinberg.
\emph{The Independence of the Continuum Hypothesis in Isabelle/ZF}.
Archive of Formal Proofs, 2022.
\url{https://isa-afp.org/entries/Independence_CH.html}.

\bibitem{afp-utm}
Jian Xu, Xingyuan Zhang, Christian Urban, Sebastiaan J. C. Joosten, and
Franz Regensburger.
\emph{Universal Turing Machine}.
Archive of Formal Proofs, 2019, revised 2022.
\url{https://isa-afp.org/entries/Universal_Turing_Machine.html}.

\bibitem{afp-models}
Emmanuel Gunther, Miguel Pagano, Pedro S\'anchez Terraf, and Mat\'ias Steinberg.
\emph{Transitive Models of Fragments of ZFC}.
Archive of Formal Proofs, 2022.
\url{https://isa-afp.org/entries/Transitive_Models.html}.

\bibitem{jech}
Thomas Jech.
\emph{Set Theory: The Third Millennium Edition, Revised and Expanded}.
Springer, 2003.
\url{https://doi.org/10.1007/3-540-44761-X}.

\bibitem{paulson-zf}
Lawrence C. Paulson.
\emph{The Relative Consistency of the Axiom of Choice, Mechanized Using
Isabelle/ZF}.
LMS Journal of Computation and Mathematics, 6, 198--248, 2003.

\bibitem{szasz}
Nora Szasz.
\emph{A Machine Checked Proof that Ackermann's Function is not Primitive
Recursive}.
In G\'erard Huet and Gordon Plotkin, editors, \emph{Logical Environments},
317--338. Cambridge University Press, 1993.

\end{thebibliography}

\end{document}
